\begin{exercisechunk}
\item \textbf{Parity orthogonality and polynomial degree.}
\label[exercise]{ex:l8-parity-orthogonality}
\exerciseprofile{2}{1}{2}{\exproof}{%
Explain why a polynomial of low degree has zero correlation with a
higher-degree parity under uniform inputs.}
Let $X$ be uniform on $\{-1,1\}^d$ and put $\chi_\varnothing(x)=1$.
\begin{enumerate}[(a),beginpenalty=10000]
\item Prove that, for all $B,C\subseteq[d]$,
\[
\E[\chi_B(X)\chi_C(X)]=\I\{B=C\}.
\]
Deduce that the $2^d$ functions $\chi_B$ form an orthonormal basis for
the vector space of real functions on $\{-1,1\}^d$, with inner product
$\langle f,g\rangle=\E[f(X)g(X)]$.

\item Let $B\ne\varnothing$, and let $p:\R^d\to\R$ be a polynomial
of total degree less than $|B|$. Prove
\[
\E[\chi_B(X)p(X)]=0.
\]
Your argument must allow repeated powers of a coordinate in a monomial.
\end{enumerate}
\begin{hints}
For (a), simplify $\chi_B(x)\chi_C(x)$.
For (b), first consider one monomial; use $X_j^2=1$.
\end{hints}
\end{exercisechunk}
